% seven = \JxFamiliesABCD (design count)
A token that states a value, such as \emph{shelf} in ``the candle is moved to the shelf'', can reach a later answer in two ways: a later token holding a candidate can match its cached key and look the value up, or the answer can copy its cached value. Clamping this writing token's key and value separately, we find that where the candidates are named decides the route: when they are listed again after the token, its key carries most of which word was written; without a later mention the identity is copied. In seven model families the route decides generated answers. Where the candidates are named also decides how deep an edit of the token can act. In three 7--8B models, overwriting the token's cached key and value from a given layer up changes free-form answers until \pctn{\NqLcFracNoneMin}--\pct{\NqLcFracNoneMax} of the depth, but answers with the candidates listed after the story only until \pctn{\NqLcFracListedMin}--\pct{\NqLcFracListedMax}; beyond that, most listed answers stay the old one, confidently, while overwriting only the layers below changes most listed answers and no free-form one. Single-site steering largely follows this window, steering at every position escapes it with a named list, and inside it an edit's share of the natural key displacement predicts whether it reaches a named-list answer. In two 7B models the matching heads write, at the later mention, a component that names no candidate and is read by key; it can move the answer but is neither specific nor the only carrier. Relative to the model's own ratio of key to value read, keys under-read steering vectors slightly and a learned remap strongly. Our analyses were preregistered or fixed in a dated note before their runs; we report estimates, and every verdict in the appendix: \JxRiskyMetABCD{} of the \JxRiskyVerdABCD{} predictions of our final preregistration judged risky before the run held (expected \JxRiskyExpABCD; \JxRiskyNotEvalABCD{} more could not be evaluated), and \NqTotSupported{} of the note's \NqTotLines{} lines were supported.
